\clearpage
\VerdierPhysicalPageAnchor{0058}%
\VerdierSemanticAnchor{ch1-1-5-10-converse}
\noindent\textbf{逆に：}\par

\VerdierSemanticAnchor{ch1-1-5-10}
\VerdierSemanticAnchor{ch1-1-5-10-proposition}
\noindent\textbf{命題 1.5.10.}\quad
\emph{基礎となる関手の間の射 $u\colon F_1\to F_2$ で、各
$i\in[1,n]$ に対して図式 (1.5.9.2) が $\beta(e_i)$-可換であるものは、
関手の $\beta$-射である。}

\VerdierSemanticAnchor{ch1-1-5-10-proof}
実際、各 $\sigma\in\mathbb Z^n$ に対し、$R(\sigma)$ で次の性質を表す：

\VerdierSemanticAnchor{ch1-1-5-10-property}
\noindent\emph{$(R(\sigma))$\qquad 次の図式は $\beta(\sigma)$-可換である：}

\VerdierSemanticAnchor{ch1-1-5-10-eq1}
\begin{equation}
  \resizebox{0.78\textwidth}{!}{%
    $\begin{tikzcd}[ampersand replacement=\&,column sep=6em,row sep=4em]
      F_1T(\sigma)
        \arrow[r,"{p_1(\sigma)}"]
        \arrow[d,"{u*T(\sigma)}"']
        \& T'(\sigma)F_1
          \arrow[d,"{T'(\sigma)*u}"] \\
      F_2T(\sigma)
        \arrow[r,"{p_2(\sigma)}"]
        \& T'(\sigma)F_2
    \end{tikzcd}$%
  }\;.
  \tag*{\ensuremath{(1.5.10.1)_{\sigma}}}
\end{equation}

図式 $(1.5.10.1)_{\sigma_1}$、$(1.5.10.1)_{\sigma_2}$、
$(1.5.10.1)_{\sigma_1+\sigma_2}$ を調べれば、直ちに次が分かる：

\VerdierSemanticAnchor{ch1-1-5-10-closure}
\noindent\emph{性質 $R(\sigma_1)$、$R(\sigma_2)$、
$R(\sigma_1+\sigma_2)$ のうち二つが成り立てば、残る一つも成り立つ。}

\VerdierSemanticAnchor{ch1-1-5-10-conclusion}
さらに仮定により、性質 $R(e_i)$（$i\in[1,n]$）および $R(0)$ が
成り立つ。したがって $R(\sigma)$ はすべての
$\sigma\in\mathbb Z^n$ に対して成り立つ。

\VerdierSemanticAnchor{ch1-1-5-11}
\noindent\textbf{1.5.11.}\quad
厳密な $\mathbb Z^n$-圏とは、その自己同型群の中の $\mathbb Z^n$ の表現を
備えた圏である。このような表現を定めるには、互いに可換する $n$ 個の自己同型
$T(e_1),\ldots,T(e_n)$ を与えれば十分である。

\VerdierSemanticAnchor{ch1-1-5-12}
\noindent\textbf{1.5.12.}\quad
与えられた圏上の（必ずしも厳密でない）$\mathbb Z^n$-圏構造を決定する問題は、
はるかに簡単ではない。例えば、圏 $\mathcal C$ 上で、すべての
$\sigma\in\mathbb Z^n$ に対して
$T(\sigma)=\text{$\mathcal C$ の恒等関手}$ となる $\mathbb Z^n$-圏構造だけを
考える。このとき各対 $\sigma_1,\sigma_2\in\mathbb Z^n$ に対して、
$\mathcal C$ の恒等関手の自己同型群（これは可換群である）の元
$c(\sigma_1,\sigma_2)$ を与え、それらが恒等式 (1.2.1.3) を満たすことを
要求しなければならない。したがって、このような構造を与えることは、
$\mathcal C$ の恒等関手の自己同型群に値をとる $\mathbb Z^n$ 上の正規化された
非斉次 2-余サイクルを与えることである。さらに、二つのこのような構造が
$F=\text{$\mathcal C$ の恒等関手}$ を満たす $\mathbb Z^n$-関手 $(F,m)$ により
同型であるための必要十分条件は、対応する余サイクルが同じコホモロジー類に
属することである。
